% 宇宙学（综述）
% license CCBYSA3
% type Wiki

本文根据 CC-BY-SA 协议转载翻译自维基百科\href{https://en.wikipedia.org/wiki/Cosmology}{相关文章}。

宇宙学（Cosmology，源自古希腊语 κόσμος（cosmos，意为“宇宙”或“世界”）与 λογία（logia，意为“研究”））是一门研究宇宙——即宇宙整体本质的物理学与形而上学的分支。英语单词 “cosmology” 首次出现在 1656 年托马斯·布朗特（Thomas Blount）的《Glossographia》中，其含义为“关于世界的论述”。\(^\text{[2]}\)1731 年，德国哲学家克里斯蒂安·沃尔夫（Christian Wolff）在拉丁语中使用 “cosmologia” 一词，用以指代研究物质世界一般性质的形而上学分支。\(^\text{[3]}\)宗教或神话宇宙学（religious or mythological cosmology）是一套基于神话、宗教与秘传文献及传统的信仰体系，通常与创世神话（creation myths）及末世论（eschatology）相关。在天文学的科学范畴中，宇宙学主要关注对宇宙年代学（chronology of the universe）的研究。

物理宇宙学（physical cosmology）研究可观测宇宙的起源、其大尺度结构与动力学特征，以及宇宙的最终命运，同时探讨支配这些领域的科学定律。\(^\text{[4]}\)这一学科由科学家（如天文学家与物理学家）以及哲学家（如形而上学家、物理哲学家与时空哲学家）共同研究。由于其研究范围与哲学存在交叉，物理宇宙学中的理论既可能包含科学命题，也可能包含非科学命题，且可能依赖于无法通过实验检验的假设。物理宇宙学是天文学的一个子分支，专注于研究整个宇宙。现代物理宇宙学由“大爆炸理论”（Big Bang Theory）主导，该理论试图结合观测天文学与粒子物理学。\(^\text{[5][6]}\)更具体而言，标准的宇宙参数化模型包括暗物质（dark matter）与暗能量（dark energy），统称为“ΛCDM 模型”（Lambda-CDM model）。

理论天体物理学家戴维·N·斯珀格尔（David N. Spergel）将宇宙学称为一门“历史科学”（historical science），因为“当我们仰望宇宙时，实际上是在回望过去”，这源于光速的有限性。\(^\text{[7]}\)
\subsection{学科领域}
\begin{figure}[ht]
\centering
\includegraphics[width=8cm]{./figures/3be309e51e6413de.png}
\caption{} \label{fig_CoGy_1}
\end{figure}
物理学与天体物理学通过科学的观测与实验，在塑造人类对宇宙的理解过程中发挥了核心作用。物理宇宙学（physical cosmology）通过数学与观测相结合的方式，在对整个宇宙的分析中逐步形成。普遍认为，宇宙起源于“大爆炸”（Big Bang），随后几乎瞬间经历了“宇宙暴涨”（cosmic inflation）——即空间的迅速膨胀。据推测，宇宙大约在 $13.799 \pm 0.021$ 十亿年前由此诞生。\(^\text{[8]}\)宇宙起源学（cosmogony）研究宇宙的起源，而宇宙描绘学（cosmography）则负责绘制宇宙的结构与特征。

在狄德罗（Diderot）的《百科全书》（Encyclopédie）中，宇宙学被划分为四个部分：天学（uranology，研究天体之科学）、气学（aerology，研究空气之科学）、地学（geology，研究大陆之科学）以及水学（hydrology，研究水体之科学）。\(^\text{[9]}\)

形而上宇宙学（metaphysical cosmology）亦被描述为探讨人类在宇宙中与其他存在物之间关系的学科。罗马皇帝兼哲学家马可·奥勒留（Marcus Aurelius）对此作出了典型的论述，他指出人的位置与认知密切相关：“不知世界为何物者，不知己所处之地；不知世界为何而存者，不知己为何人，亦不知世界为何物。”\(^\text{[10]}\)
\subsection{发现}
\subsubsection{物理宇宙学}
物理宇宙学是物理学与天体物理学的一个分支，研究宇宙的物理起源与演化，同时也探讨宇宙在大尺度上的性质。在其最早期的形式中，它即如今所谓的“天体力学”（celestial mechanics）——对天体运动与“天界”的研究。希腊哲学家萨摩斯的阿里斯塔克（Aristarchus of Samos）、亚里士多德（Aristotle）与托勒密（Ptolemy）提出了不同的宇宙理论。其中，托勒密的地心体系（geocentric Ptolemaic system）长期占据主导地位，直到十六世纪，尼古拉·哥白尼（Nicolaus Copernicus）以及随后约翰内斯·开普勒（Johannes Kepler）与伽利略·伽利莱（Galileo Galilei）提出日心体系（heliocentric system）。这是物理宇宙学中最著名的认识论断裂（epistemological rupture）之一。

艾萨克·牛顿（Isaac Newton）于 1687 年发表的《自然哲学的数学原理》（Philosophiæ Naturalis Principia Mathematica）首次系统描述了万有引力定律（law of universal gravitation）。该理论为开普勒行星运动定律提供了物理机制，并能够解释行星之间引力相互作用所导致的系统异常。牛顿宇宙学与其前辈理论的一个根本区别在于哥白尼原理（Copernican principle）：地球上的物体遵循与天体相同的物理定律。这一思想构成了物理宇宙学的一个重要哲学飞跃。

现代科学宇宙学通常被认为始于 1917 年，当时阿尔伯特·爱因斯坦（Albert Einstein）发表了其对广义相对论（general relativity）的最终修正论文《广义相对论的宇宙学考量》（Cosmological Considerations of the General Theory of Relativity）。\(^\text{[11]}\)（尽管这篇论文直到第一次世界大战结束后才在德国以外被广泛传播。）广义相对论的提出促使威廉·德西特（Willem de Sitter）、卡尔·史瓦西（Karl Schwarzschild）与阿瑟·爱丁顿（Arthur Eddington）等宇宙起源学家（cosmogonists）探讨其在天文学中的影响，从而极大地提升了天文学家研究遥远天体的能力。物理学家们也由此开始放弃“宇宙是静止且不变的”这一传统假设。1922 年，亚历山大·弗里德曼（Alexander Friedmann）首次提出“宇宙膨胀”概念，即宇宙中包含运动的物质并在不断扩张。

与这种动态宇宙学观点并行的，是关于宇宙结构的一场长期争论——即 1917 年至 1922 年间的“大辩论”（Great Debate）。在这一时期，早期的宇宙学家如赫伯·柯蒂斯（Heber Curtis）与恩斯特·厄皮克（Ernst Öpik）通过观测确定，望远镜中所见的一些星云（nebulae）其实是距离我们极其遥远的独立星系。\(^\text{[12]}\) 柯蒂斯主张旋涡星云（spiral nebulae）本身就是独立的恒星系统，即“岛宇宙”（island universes）；而威尔逊山天文台的天文学家哈洛·沙普利（Harlow Shapley）则坚持认为宇宙仅由银河系这一恒星系统构成。两种观点的分歧在 1920 年 4 月 26 日于华盛顿特区美国国家科学院（U.S. National Academy of Sciences）会议上举行的“大辩论”中达到高潮。此争论最终在 1923 至 1924 年间得到解决，当时埃德温·哈勃（Edwin Hubble）在仙女座星系（Andromeda Galaxy）中探测到造父变星（Cepheid Variables）。\(^\text{[13][14]}\)其测得的距离证实了旋涡星云远远位于银河系边界之外。

此后，对宇宙的理论建模进一步探索了这样一种可能性：爱因斯坦在其 1917 年论文中引入的宇宙常数（cosmological constant），其数值可能导致宇宙处于膨胀状态。由此，比利时神父乔治·勒梅特（Georges Lemaître）于 1927 年提出了“大爆炸模型”（Big Bang model）。\(^\text{[15]}\)这一理论随后得到了哈勃 1929 年发现红移现象（redshift）的支持，\(^\text{[16]}\)并在 1964 年由阿诺·彭齐亚斯（Arno Penzias）与罗伯特·威尔逊（Robert Woodrow Wilson）发现宇宙微波背景辐射（cosmic microwave background radiation）后得到进一步验证。\(^\text{[17]}\)这些发现首次排除了多种其他竞争性的宇宙学模型。

自 1990 年前后起，观测宇宙学（observational cosmology）取得了多项重大突破，使宇宙学从主要依赖推测的学科转变为能够进行高精度预测的科学领域——理论与观测之间达成了前所未有的精确一致。这些进展包括：来自 COBE、\(^\text{[18]}\)WMAP、\(^\text{[19]}\)与 Planck\(^\text{[20]}\)卫星的宇宙微波背景观测；大型星系红移巡天（galaxy redshift surveys），如 2dFGRS\(^\text{[21]}\)与 SDSS\(^\text{[22]}\)；以及对遥远超新星与引力透镜（gravitational lensing）的观测。这些观测结果与“宇宙暴胀理论”（cosmic inflation theory）——即修正的大爆炸理论及其特定版本“ΛCDM 模型”——的预测高度吻合。因此，许多学者称当代为“宇宙学的黄金时代”（golden age of cosmology）。\(^\text{[23]}\)

2014 年，BICEP2 合作组宣称他们在宇宙微波背景中探测到了引力波（gravitational waves）的印记。然而，后续研究表明该结果并不可靠：所谓的引力波证据实际上是由星际尘埃造成的。\(^\text{[24][25]}\)

同年 12 月 1 日，在意大利费拉拉（Ferrara）召开的 Planck 2014 会议上，天文学家宣布：宇宙的年龄约为 138 亿年，其组成包括 4.9\% 的原子物质、26.6\% 的暗物质（dark matter），以及 68.5\% 的暗能量（dark energy）。\(^\text{[26]}\)
\subsubsection{宗教或神话宇宙学}
宗教或神话宇宙学是一套基于神话、宗教与秘传文献以及创世与末世传统（creation and eschatology）的信仰体系。几乎所有宗教都包含创世神话（creation myths），并且通常可依据米尔恰·伊利亚德（Mircea Eliade）及其同事查尔斯·朗（Charles Long）所建立的分类体系分为五种类型。

基于相似主题的创世神话类型（Types of Creation Myths based on similar motifs）包括：
\begin{enumerate}
    \item \textbf{无中生有的创造（Creation ex nihilo）}：创造源自神祇的思想、言语、梦境或身体分泌物。
    \item \textbf{潜水取地型创造（Earth diver creation）}：创造者派出潜水者（通常是鸟类或两栖动物）潜入原始之海的海底，从中取出沙土或泥浆，形成陆地世界。
    \item \textbf{逐层显现型创造（Emergence myths）}：始祖（progenitors）经历多个世界与形态的转变，直至抵达现今的世界。
    \item \textbf{肢解型创造（Creation by dismemberment）}：世界由原初存在（primordial being）的肢解与分裂而形成。
    \item \textbf{原初统一体分裂型创造（Creation by splitting or ordering of a primordial unity）}：宇宙由原始统一体（如“宇宙卵”的破裂或从混沌中建立秩序）而生。\(^\text{[27]}\)
\end{enumerate}
\subsubsection{哲学}
\begin{figure}[ht]
\centering
\includegraphics[width=6cm]{./figures/9ad29988d4f7e015.png}
\caption{以对数尺度表示的可观测宇宙示意图（Representation of the observable universe on a logarithmic scale）。从中心向外，距离由太阳开始逐渐增加。行星及其他天体的尺寸被放大，以便更清晰地展示其形状。} \label{fig_CoGy_2}
\end{figure}
宇宙学（Cosmology）研究作为整体的世界——包括空间、时间及一切现象。从历史上看，其研究范围极为广泛，并在许多情形下与宗教密切相关。\(^\text{[28]}\)某些关于宇宙的问题超出了科学探究的范围，但仍可通过其他哲学方法（如辩证法 dialectics）加以思考。超出科学范畴的探讨通常包括如下问题：\(^\text{[29][30]}\)
\begin{itemize}
    \item 宇宙的起源是什么？其“第一因”是否存在？其存在是否必然？（参见：一元论 monism、泛神论 pantheism、流溢论 emanationism、创世论 creationism）
    \item 宇宙的最终物质成分是什么？（参见：机械论 mechanism、动力论 dynamism、质形论 hylomorphism、原子论 atomism）
    \item 宇宙存在的最终理由（若有）是什么？宇宙是否具有目的？（参见：目的论 teleology）
    \item 意识的存在在现实的存在中是否起作用？我们如何得知自己关于宇宙整体的知识？宇宙学推理是否揭示形而上学真理？（参见：认识论 epistemology）
\end{itemize}
哲学史学者查尔斯·卡恩（Charles Kahn）将古希腊宇宙学的起源归因于阿那克西曼德（Anaximander）。\(^\text{[31]}\)





\subsection{历史宇宙学}
\begin{table}[ht]
\centering
\caption{}\label{tab_CoGy1}
\begin{tabular}{|c|c|c|c|}
\hline
\textbf{名称} & \textbf{作者与年代} & \textbf{分类} & \textbf{备注} \\
\hline
印度宇宙学 & 《梨俱吠陀》（Rigveda，约公元前 1700–1100 年）& 循环或振荡型（Cyclical or oscillating），时间上无限（Infinite in time） & 原初物质显现 311.04 万亿年后进入同等时长的非显现状态。宇宙显现 43.2 亿年后亦进入同等时长的非显现状态。无数个宇宙同时存在，这些循环自始至终永恒持续，由“欲望”（desires）所驱动。 \\
\hline
祆教宇宙学（Zoroastrian Cosmology）&《阿维斯陀经》（Avesta，约公元前 1500–600 年）& 二元宇宙论（Dualistic Cosmology） & 根据祆教宇宙学，宇宙是“存在与非存在”、“善与恶”、“光与暗”之间永恒冲突的显现。宇宙将在这种对立状态中持续 12000 年；在复活之时，这两种元素将再次被分离。\\
\hline
耆那教宇宙学（Jain cosmology）&《耆那教经典》（Jain Agamas，约公元 500 年撰写，依据大雄 Mahavira〔公元前 599–527 年〕的教义） & 循环或振荡型（Cyclical or oscillating），永恒而有限（eternal and finite） & 耆那教宇宙学认为“娑罗迦”（loka，即宇宙）是一个未被创造的永恒存在，其形状类似一个双腿分开、手臂支于腰间的站立之人。根据耆那教的描述，宇宙的上部宽阔，中部狭窄，而底部再次变得宽广。\\
\hline
巴比伦宇宙学（Babylonian cosmology）& 巴比伦文献（Babylonian literature，约公元前 2300–500 年）& 浮于无限“混沌之水”上的平面地球（Flat Earth floating in infinite "waters of chaos"）& 地球与天穹共同构成一个位于无限“混沌之水”（waters of chaos）中的整体；地球为平面且呈圆形，其上方覆盖着坚固的穹顶（即“天穹” firmament），以阻隔外部的“混沌之海”。\\
\hline
埃利亚学派宇宙学（Eleatic cosmology）& 巴门尼德（Parmenides，约公元前 515 年）  & 有限且球形（Finite and spherical in extent）& 宇宙是不变的、均匀的、完美的、必然的、超越时间的，既非生成也非消逝；虚空（Void）不存在。多样性与变化皆源于感官经验所导致的认知无知（epistemic ignorance）。时间与空间的界限皆为任意设定，且相对于巴门尼德式的整体（Parmenidean whole）而言是相对的。\\
\hline
数论派宇宙演化论（Samkhya Cosmic Evolution）& 迦毗罗（Kapila，公元前 6 世纪），阿修里（Asuri）之弟子 & “自性”（Prakriti，物质）与“神我”（Purusha，意识）之关系（Relation between Matter and Consciousness）& “自性”（Prakriti，物质）是生成世界（world of becoming）的根源，是一种纯粹的潜能（pure potentiality），依次演化为二十四种“原理”（tattvas）。这种演化之所以可能，是因为 Prakriti 始终处于其三种组成要素（称为“性” gunas）之间的张力状态之中：善性（Sattva，轻明或纯净）、动性（Rajas，激情或活动）、惰性（Tamas，惯性或沉重）。数论派的因果理论被称为“有因论”（Satkārya-vāda，theory of existent causes），其核心主张为：无物可凭空生灭，所有的演化不过是原始自然（primal Nature）从一种形式向另一种形式的转化。\\
\hline
圣经宇宙学（Biblical cosmology）&《创世纪》创世叙事（Genesis creation narrative）& 漂浮于无限“混沌之水”中的地球（Earth floating in infinite "waters of chaos"）& 地球与天穹构成一个位于无限“混沌之水”中的整体；其上方的“穹苍”（firmament）阻隔外部的“混沌之海”。\\
\hline
阿那克西曼德模型（Anaximander's model）& 阿那克西曼德（Anaximander，约公元前 560 年）& 地心说（Geocentric），圆柱形地球，空间无限而时间有限；首个纯机械模型（first purely mechanical model）& 地球静止地漂浮于无限空间的中心，不受任何支撑。\(^\text{[32]}\)起初，在冷热分离之后，一团火球出现，像树皮包裹树干一样环绕着地球。此火球破裂后形成了宇宙的其他部分。宇宙如同一个由火充盈的空心同心圆环系统，其边缘布满洞孔，如笛孔一般；并无真正的天体，唯有火光透过孔隙射出。共有三层圆环，从地球向外依次为：恒星（包括行星）、月亮与巨大的太阳。\(^\text{[33]}\)\\
\hline
原子论宇宙（Atomist universe）& 阿那克萨哥拉（Anaxagoras，公元前 500–428 年）及后来的伊壁鸠鲁（Epicurus）& 无限宇宙（Infinite in extent）& 宇宙中仅包含两种存在：无限数量的微小“种子”（即原子 atoms）与无限广阔的虚空（void）。所有原子由同一物质构成，但在大小与形状上有所差异。万物由原子的聚合而生，最终又分解为原子。此理论吸收了留基伯（Leucippus）的因果原理：“无事偶然，一切因理与必然而生。” 宇宙不受诸神统治。\(^\text{[citation needed]}\)\\
\hline
毕达哥拉斯宇宙（Pythagorean universe）& 腓罗劳斯（Philolaus，卒于公元前 390 年）& 存在“中央之火”（Central Fire）于宇宙中心（Existence of a “Central Fire” at the center of the Universe）& 在宇宙的中心存在一团“中央之火”，地球、太阳、月亮与行星皆以匀速围绕其旋转。太阳每年绕中央之火一周，而恒星则静止不动。地球在运行中始终以同一面朝向中央之火，因此该火永不为人所见。这是已知最早的“非地心宇宙模型”。\(^\text{[34]}\)\\
\hline
《论世界》（De Mundo）& 伪亚里士多德（Pseudo-Aristotle，卒于公元前 250 年，或约公元前 350–200 年间）& 宇宙是由天与地及其中的元素构成的系统（The Universe is a system made up of heaven and Earth and the elements contained within them）& 宇宙由“五种元素”组成，分布于五个球层中：地（earth）被水（water）包围，水被气（air）包围，气被火（fire）包围，而火又被以太（ether）所包围，这五层共同构成整个宇宙。\(^\text{[35]}\)\\
\hline
斯多亚派宇宙（Stoic universe）& 斯多亚学派（Stoics，公元前 300 年至公元 200 年）& 岛宇宙（Island universe）& 宇宙是有限的，被无限虚空（infinite void）所环绕。宇宙处于不断流变之中，其规模脉动变化，周期性地经历剧烈的动荡与宇宙火劫（conflagrations）。\\
\hline
柏拉图宇宙（Platonic universe）& 柏拉图（Plato，约公元前 360 年）& 地心说（Geocentric），复杂宇宙生成论（complex cosmogony），空间有限（finite extent），暗含有限时间（implied finite time），循环性（cyclical）& 静止的地球位于宇宙中心，其外依次环绕天体，这些天体在完美的圆周上运动，由造物主（Demiurge）依意志所安排：月亮、太阳、行星与恒星。\(^\text{[36][37][38]}\)复杂的天体运动在每一个“完美之年”中循环重复。\(^\text{[39]}\)\\
\hline
欧多克索斯模型（Eudoxus' model）& 克尼多斯的欧多克索斯（Eudoxus of Cnidus，约公元前 340 年）及后来的卡利普斯（Callippus）& 地心说（Geocentric），首个几何—数学模型（first geometric-mathematical model）& 天体的运动如同附着在多个以地球为中心的同心不可见球体上，每个球体都围绕自身不同的轴线、以不同的速度旋转。\(^\text{[40]}\)模型共设有二十七个同心球（homocentric spheres），每个球体用以解释特定天体的一类可观测运动。欧多克索斯强调这些球体纯为数学构造——它们在物理上并不存在，仅用于表征天体可能出现的位置。\(^\text{[41]}\)\\
\hline
亚里士多德宇宙（Aristotelian universe）& 亚里士多德（Aristotle，公元前 384–322 年）& 地心说（基于欧多克索斯模型），静态（static）、稳态（steady state）、空间有限（finite extent）、时间无限（infinite time）& 静止且呈球形的地球被 43 至 55 层同心天球所包围，这些天球为物质性的水晶结构。\(^\text{[42]}\)宇宙自永恒以来保持不变。其体系包含第五元素——“以太”（aether），作为对四种经典元素的补充。\(^\text{[43]}\)\\
\hline
阿里斯塔克斯宇宙（Aristarchean universe）& 阿里斯塔克斯（Aristarchus，约公元前 280 年）& 日心说（Heliocentric）& 地球每日自转一周，并以圆形轨道每年绕太阳公转一次。恒星球（sphere of fixed stars）以太阳为中心。\(^\text{[44]}\)\\
\hline
托勒密模型（Ptolemaic model）& 托勒密（Ptolemy，公元 2 世纪）& 地心说（Geocentric，基于亚里士多德宇宙）& 宇宙围绕静止的地球运转。行星在称为“本轮”（epicycle）的圆轨道上运动，而每个本轮的圆心又沿着更大的圆轨道（称为偏轮 eccentric 或均轮 deferent）绕地球附近的中心点运行。引入“均差点”（equants）的概念进一步提高了模型的复杂度，使天文学家得以精确预测行星位置。此模型因其长期稳定的预测能力，被认为是史上最成功的宇宙模型之一，其主要著作为《天文学大成》（Almagest，又称“大体系” The Great System）。\\
\hline
卡佩拉模型（Capella's model）& 马提安努斯·卡佩拉（Martianus Capella，约公元 420 年）& 地心与日心混合模型（Geocentric and Heliocentric）& 地球静止于宇宙中心，月亮、太阳、三颗行星及恒星环绕其运行，而水星与金星则绕太阳公转。\(^\text{[45]}\)\\
\hline
阿耶波多模型（Aryabhatan model）& 阿耶波多（Aryabhata，公元 499 年）& 地心或日心模型（Geocentric or Heliocentric）& 地球自转，行星以椭圆轨道绕地球或太阳运行；由于行星轨道的位置既以地球又以太阳为参照，模型究竟是地心还是日心尚不确定。\\
\hline
《古兰经》宇宙学（Quranic cosmology）& 《古兰经》（Quran，公元 610–632 年）& 平面地球观（Flat-earth）& 宇宙由叠层的平面结构构成，其中包括七重天（seven levels of heaven），在部分诠释中还包含七层地（seven levels of earth），其中最下层为地狱（hell）。\\
\hline
中世纪宇宙（Medieval universe）& 中世纪哲学家（公元 500–1200 年） &时间上有限（Finite in time）& 基督教学者约翰·菲罗波诺斯（John Philoponus）提出宇宙在时间上是有限的并具有起点，从而反对古希腊“无限过去”的概念。早期穆斯林哲学家阿尔·肯迪（Al-Kindi）、犹太哲学家萨阿迪亚·高恩（Saadia Gaon）以及穆斯林神学家阿尔·加扎利（Al-Ghazali）进一步发展了支持有限宇宙的逻辑论证。\\
\hline
非平行多元宇宙（Non-Parallel Multiverse）& 《博伽梵往世书》（Bhagvata Purana，公元 800–1000 年）& 多元宇宙，非平行结构（Multiverse, Non-Parallel）& 无数宇宙的存在可与现代多元宇宙理论相比较，但这些宇宙并非平行——每个宇宙皆独立不同，每一具个体灵魂（jīva-ātma，有身之魂）在任一时刻仅存在于一个宇宙之中。所有宇宙均由同一物质生成，因此它们遵循平行的时间循环——同时显现与隐没。\(^\text{[46]}\)\\
\hline
多宇宙宇宙学（Multiversal cosmology）& 法赫尔·丁·拉齐（Fakhr al-Din al-Razi，1149–1209）& 多元宇宙、多世界（Multiverse, multiple worlds and universes）& 在已知世界之外存在无限外层空间，真主（God）拥有能力使真空中充满无数个宇宙。\\
\hline
马拉噶模型（Maragha models）& 马拉噶天文学派（Maragha school，1259–1528）& 地心说（Geocentric）& 对托勒密模型（Ptolemaic model）与亚里士多德宇宙（Aristotelian universe）进行了多项修正，包括在马拉噶天文台（Maragheh observatory）中弃用均差点（equant）与偏心圆（eccentrics），并由纳西尔·丁·图西（Al-Tusi）引入“图西双圆”（Tusi-couple）。其后又提出若干替代模型，如伊本·沙提尔（Ibn al-Shatir）提出的首个精确月球运动模型，阿里·库什丘（Ali Kuşçu）提出以地球自转取代静止地球的理论，以及阿尔·比尔詹迪（Al-Birjandi）提出包含“圆周惯性”（circular inertia）概念的行星模型。\\
\hline
尼拉坎塔模型（Nilakanthan model）& 尼拉坎塔·索马亚吉（Nilakantha Somayaji，1444–1544）& 地心与日心混合模型（Geocentric and heliocentric）& 行星绕太阳公转，而太阳再绕地球运行的宇宙体系，与后来的第谷体系（Tychonic system）相似。\\
\hline
哥白尼宇宙（Copernican universe）& 尼古拉·哥白尼（Nicolaus Copernicus，1473–1543）& 日心说，行星圆轨道，空间有限（Heliocentric with circular planetary orbits, finite extent）& 首次发表于《天体运行论》（De revolutionibus orbium coelestium）。太阳位于宇宙中心，地球及其他行星绕太阳公转，而月球绕地球运动。宇宙的外边界由“恒星天球”（sphere of fixed stars）限定。\\
\hline
第谷体系（Tychonic system）& 第谷·布拉赫（Tycho Brahe，1546–1601）& 地心与日心混合模型（Geocentric and heliocentric）& 行星绕太阳公转，而太阳再绕地球运行，与早期的尼拉坎塔模型相似。\\
\hline
布鲁诺宇宙学（Bruno's cosmology）& 焦尔达诺·布鲁诺（Giordano Bruno，1548–1600）& 空间无限、时间无限、均质、各向同性、无层级（Infinite extent, infinite time, homogeneous, isotropic, non-hierarchical）& 否定宇宙的层级结构，认为地球与太阳与其他天体相比并无特殊之处。恒星之间的虚空充满以太（aether），物质由同样的四种元素——水、土、火、气——组成，并具备原子性（atomistic）、泛灵性（animistic）与理性（intelligent）。\\
\hline
《论磁体》（De Magnete）& 威廉·吉尔伯特（William Gilbert，1544–1603）& 日心说，空间无限延展（Heliocentric, indefinitely extended）& 继承哥白尼的日心体系，但拒绝“恒星天球”作为宇宙边界的假设，认为无证据支持其存在。\(^\text{[47]}\)\\
\hline
开普勒体系（Keplerian）& 约翰内斯·开普勒（Johannes Kepler，1571–1630）& 日心说，行星椭圆轨道（Heliocentric with elliptical planetary orbits）& 开普勒的发现将数学与物理学相结合，为现代太阳系结构奠定基础，但当时的远方恒星仍被视为附着于薄且静止的天球之上。\\
\hline
静态牛顿宇宙（Static Newtonian）& 艾萨克·牛顿（Isaac Newton，1642–1727）& 静态（随时间演化）、稳态（steady state）、无限（infinite）& 宇宙中每个粒子都吸引其他所有粒子。物质在大尺度上均匀分布。体系在引力上达到平衡，但这种平衡是不稳定的。\\
\hline
笛卡尔漩涡宇宙（Cartesian Vortex universe）& 勒内·笛卡尔（René Descartes，17 世纪）& 静态（随时间演化）、稳态、无限 & 宇宙由巨大的以太或精微物质（aethereal or fine matter）漩涡体系构成，这些旋转漩涡产生引力效应。然而其真空并非空无一物，整个空间都充满物质。\\
\hline
层级宇宙（Hierarchical universe）& 伊曼努尔·康德（Immanuel Kant）、约翰·兰伯特（Johann Lambert），18 世纪 & 静态（随时间演化）、稳态、无限 & 物质在不断扩大的层级结构中聚集。物质在各层之间无尽地循环再生。\\
\hline
爱因斯坦宇宙（带宇宙常数）（Einstein Universe with a cosmological constant）& 阿尔伯特·爱因斯坦（Albert Einstein，1917）& 名义上静态（Static, nominally），有限（Bounded, finite）& “无运动的物质”（Matter without motion）。包含均匀分布的物质。空间为均匀弯曲的球形空间，基于黎曼的高维球体（Riemann hypersphere）。曲率由宇宙常数 $\Lambda$ 决定；实际上 $\Lambda$ 等效于一种抵消引力的排斥力。该模型不稳定。\\
\hline
德西特宇宙（De Sitter universe）& 威廉·德西特（Willem de Sitter，1917）& 膨胀的平坦空间（Expanding flat space），稳态，$\Lambda > 0$ & “无物质的运动”（Motion without matter）。仅表面上静态。基于爱因斯坦广义相对论（general relativity）。空间以恒定加速度膨胀，尺度因子呈指数增长（持续暴胀）。\\
\hline
麦克米伦宇宙（MacMillan universe）& 威廉·邓肯·麦克米伦（William Duncan MacMillan，1920 年代）& 静态且稳态（Static and steady state）& 新物质由辐射（radiation）转化而成；恒星光不断循环，被再生为新的物质粒子。\\
\hline
弗里德曼宇宙（球形空间）（Friedmann universe, spherical space） & 亚历山大·弗里德曼（Alexander Friedmann，1922）& 球形膨胀空间（Spherical expanding space），$k = +1$，无宇宙常数 $\Lambda$ & 正曲率，曲率常数 $k = +1$。宇宙先膨胀后收缩。空间封闭（有限）。\\
\hline
弗里德曼宇宙（双曲空间）（Friedmann universe, hyperbolic space） & 亚历山大·弗里德曼（Alexander Friedmann，1924）& 双曲膨胀空间（Hyperbolic expanding space），$k = -1$，无宇宙常数 $\Lambda$ & 负曲率，据称是无限的（但解释上含糊），空间无边界，将永远膨胀。\\
\hline
狄拉克大数假说（Dirac large numbers hypothesis）& 保罗·狄拉克（Paul Dirac，1930 年代）& 膨胀宇宙（Expanding）& 该理论要求引力常数 $G$ 随时间显著变化，并随着宇宙演化逐渐减小，即引力随时间推移而减弱。\\
\hline
弗里德曼零曲率宇宙（Friedmann zero-curvature）& 爱因斯坦与德西特（Einstein and De Sitter，1932）& 膨胀的平坦空间（Expanding flat space），$k = 0$，$\Lambda = 0$，临界密度（Critical density）& 曲率常数 $k = 0$。据称该宇宙是无限的（但解释上含糊），被描述为“无界而有限的宇宙”（Unbounded cosmos of limited extent）。该模型永远膨胀，是所有已知宇宙模型中“最简单”的一种。以弗里德曼命名，但非其本人提出。其减速因子 $q = 1/2$，意味着宇宙的膨胀速率在逐渐减慢。\\
\hline
原始大爆炸模型（The original Big Bang, Friedmann–Lemaître）& 乔治·勒梅特（Georges Lemaître，1927–1929）& 膨胀宇宙（Expansion），$\Lambda > 0$，且 $\Lambda > |\text{Gravity}|$ & 宇宙常数 $\Lambda$ 为正，且其量级大于引力。宇宙从一个高密度初态（“原初原子”primeval atom）起始，随后经历两个阶段的膨胀。$\Lambda$ 被用于“使宇宙失稳”以促成膨胀。（勒梅特被公认为“大爆炸模型之父”。）\\
\hline
振荡宇宙（Friedmann–Einstein Oscillating universe）& 弗里德曼主张于 1920 年代 & 膨胀与收缩交替循环（Expanding and contracting in cycles）& 时间无始无终，从而避免“时间起点悖论”。宇宙经历永恒的循环：大爆炸（Big Bang）之后是大坍缩（Big Crunch）。（这是爱因斯坦在否定其 1917 年静态模型后提出的首选方案。）\\
\hline
爱丁顿宇宙（Eddington universe）& 阿瑟·爱丁顿（Arthur Eddington，1930） & 先静态后膨胀（First static then expands）& 基于爱因斯坦 1917 年的不稳定静态宇宙，其平衡被扰动后进入膨胀状态；随着物质的不断稀释，演化为德西特宇宙。宇宙常数 $\Lambda$ 支配引力效应。\\
\hline
米尔恩运动相对论宇宙（Milne universe of kinematic relativity）& 爱德华·米尔恩（Edward Milne，1933, 1935）；威廉·H·麦克雷（William H. McCrea，1930s） & 运动学膨胀（Kinematic expansion without space expansion）& 否定广义相对论与空间膨胀范式。模型未将引力作为初始假设。符合宇宙学原理（cosmological principle）与狭义相对论（special relativity）。宇宙由一个有限的球形粒子（或星系）云构成，在无限且空旷的平坦空间中膨胀。它具有中心与边界（粒子云表面），边界以光速膨胀。其引力解释复杂且缺乏说服力。\\
\hline
弗里德曼–勒梅特–罗伯逊–沃尔克模型族（Friedmann–Lemaître–Robertson–Walker class of models）& 霍华德·罗伯逊（Howard Robertson）、阿瑟·沃尔克（Arthur Walker，1935）& 均匀膨胀宇宙（Uniformly expanding）& 一类均匀且各向同性的宇宙模型。时空分解为均匀曲率空间与宇宙时间，所有共动观察者共享同一宇宙时。此系统化形式被称为弗里德曼–勒梅特–罗伯逊–沃尔克度规（FLRW metrics）或罗伯逊–沃尔克度规。\\
\hline
稳态宇宙（Steady-state, Bondi–Gold）& 赫尔曼·邦迪（Hermann Bondi）、托马斯·戈尔德（Thomas Gold，1948）& 膨胀、稳态、无限（Expanding, steady state, infinite）& 物质的生成速率维持恒定密度。物质不断从“无中”产生（continuous creation out of nothing from nowhere）。宇宙呈指数膨胀。减速因子 $q = -1$。\\
\hline
稳态宇宙（Steady-state, Hoyle）& 弗雷德·霍伊尔（Fred Hoyle，1948）& 膨胀、稳态但不稳定（Expanding, steady state; but unstable）& 物质的生成速率维持恒定密度。然而，由于物质生成速率必须与空间膨胀速率精确平衡，该体系在本质上是不稳定的。\\
\hline
双重等离子宇宙（Ambiplasma universe）& 汉尼斯·阿尔文（Hannes Alfvén）、奥斯卡·克莱因（Oskar Klein，1965）& 等离子体宇宙（Cellular universe, expanding by matter–antimatter annihilation）& 基于等离子体宇宙学（plasma cosmology）。宇宙被视为由“双层结构”分隔的“超星系”（meta-galaxies）组成的泡状体系。其他宇宙由其他泡泡形成。持续的物质–反物质湮灭（matter–antimatter annihilation）使这些泡泡保持分离并不断远离，从而防止它们相互作用。\\
\hline
布兰斯–迪克理论（Brans–Dicke theory）& 卡尔·H·布兰斯（Carl H. Brans）、罗伯特·H·迪克（Robert H. Dicke）& 膨胀宇宙（Expanding） & 基于马赫原理（Mach's principle）。引力常数 $G$ 随宇宙膨胀而变化。然而，“没有人能确切说明马赫原理究竟意味着什么。”\\
\hline
宇宙暴胀（Cosmic inflation）& 阿兰·古思（Alan Guth，1980）& 修改后的大爆炸模型，用以解决视界问题与平坦性问题（Big Bang modified to solve horizon and flatness problems）& 基于“热暴胀”（hot inflation）的概念。宇宙被视为多重量子涨落（multiple quantum flux）的结果，因此呈现泡状结构（bubble-like nature）。其他宇宙由其他泡泡形成。持续的宇宙膨胀使这些泡泡保持分离并相互远离。\\
\hline
永恒暴胀（Eternal inflation，亦即多元宇宙模型）& 安德烈·林德（Andreï Linde，1983）& 含宇宙暴胀的大爆炸模型（Big Bang with cosmic inflation）& 多元宇宙（Multiverse）基于“冷暴胀”（cold inflation）的概念，其中暴胀事件随机发生，每次具有独立的初始条件；部分泡泡膨胀成类似整个宇宙的“泡泡宇宙”（bubble universes）。这些泡泡在时空泡沫（spacetime foam）中不断成核（nucleate）。\\
\hline
循环宇宙模型（Cyclic model）& 保罗·斯坦哈特（Paul Steinhardt）、尼尔·图罗克（Neil Turok，2002）& 膨胀与收缩周期交替；基于 M 理论（Expanding and contracting in cycles; M-theory）& 两个平行的轨道平面或 M-膜（M-branes）在高维空间中周期性碰撞。模型包含“第五元素能”（quintessence）或暗能量（dark energy）。\\
\hline
循环宇宙模型（Cyclic model）& 劳里斯·鲍姆（Lauris Baum）、保罗·弗兰普顿（Paul Frampton，2007）& 托尔曼熵问题的解（Solution of Tolman’s entropy problem）& 幻影暗能量（phantom dark energy）使宇宙分裂为大量彼此不相连的区域。可观测区域在收缩过程中仅保留暗能量且熵为零。\\
\hline
\end{tabular}
\end{table}
表格注释：术语“静态”（static）仅表示宇宙既不膨胀也不收缩。符号 $G$ 表示牛顿引力常数（Newton's gravitational constant）；$\Lambda$（Lambda）为宇宙常数（cosmological constant）。
\subsection{另见 } 
\begin{itemize}
\item 绝对时间与空间（Absolute time and space）  
\item 大历史（Big History）  
\item 地球科学（Earth science）  
\item 星系形成与演化（Galaxy formation and evolution）  
\item Illustris 项目（Illustris project）  
\item 耆那教与非创世论（Jainism and non-creationism）  
\item $\Lambda$–CDM 模型（Lambda–CDM model）  
\item 天体物理学家列表（List of astrophysicists）  
\item 非标准宇宙学（Non-standard cosmology）  
\item 太极哲学（Taiji (philosophy)）  
\item 宇宙学理论时间线（Timeline of cosmological theories）  
\item 普适旋转曲线（Universal rotation curve）  
\item 暖暴胀理论（Warm inflation）  
\item “大环”结构（Big Ring）
\end{itemize}
\subsection{参考文献}
\begin{enumerate}
\item Hille, Karl, ed. (2016 年 10 月 13 日). “哈勃揭示可观测宇宙中星系数量是先前估计的十倍。”（"Hubble Reveals Observable Universe Contains 10 Times More Galaxies Than Previously Thought"）NASA. 取自 2016 年 10 月 17 日。  
\item Hetherington, Norriss S. (2014). 《宇宙学百科全书（Routledge 重刊版）：现代宇宙学的历史、哲学与科学基础》（Encyclopedia of Cosmology (Routledge Revivals): Historical, Philosophical, and Scientific Foundations of Modern Cosmology）. Routledge. 第 116 页. ISBN 978-1-317-67766-6。  
\item Luminet, Jean-Pierre (2008). 《包裹的宇宙》（The Wraparound Universe）. CRC Press. 第 170 页. ISBN 978-1-4398-6496-8。  
\item “导论：宇宙学——空间”（"Introduction: Cosmology – space"）《新科学人》（New Scientist）, 2006 年 9 月 4 日。网页存档于 2015 年 7 月 3 日。  
\item “Cosmology.” 《牛津词典》（Oxford Dictionaries）。  

\item Overbye, Dennis (2019 年 2 月 25 日). “暗能量是否正在扰乱宇宙？——轴子？幻影能量？天体物理学家试图修补宇宙的漏洞，并在此过程中重写宇宙史。”《纽约时报》（*The New York Times*）。  

\item Spergel, David N. (2014 秋季). “今日宇宙学”（"Cosmology Today"）. *Daedalus*, 143 (4): 125–133. doi:10.1162/DAED_a_00312。  

Planck Collaboration (2016 年 10 月 1 日). “普朗克 2015 结果. XIII. 宇宙学参数”（"Planck 2015 results. XIII. Cosmological parameters"）. *Astronomy & Astrophysics*, 594 (13): 表 4, PDF 第 31 页。arXiv:1502.01589。  

Diderot, Denis (2015 年 4 月 1 日). “人类知识体系的详细解释”（"Detailed Explanation of the System of Human Knowledge"）. *Encyclopedia of Diderot & d'Alembert – Collaborative Translation Project*.  

《马可·奥勒留沉思录》第八卷第 52 节（*The Thoughts of Marcus Aurelius Antoninus viii.52*）。  

Einstein, Albert (1952). “论广义相对论的宇宙学考虑”（"Cosmological considerations on the general theory of relativity"）. 收录于 *The Principle of Relativity*. Dover. 第 175–188 页。  

Dodelson, Scott (2003 年 3 月 30 日). 《现代宇宙学》（*Modern Cosmology*）. Elsevier. ISBN 978-0-08-051197-9。  

Falk, Dan (2009 年 3 月 18 日). “书评：《我们发现宇宙的那一天》作者 Marcia Bartusiak。”（"Review: The Day We Found the Universe"）. *New Scientist*, 201 (2700): 45。  

Hubble, E. P. (1926 年 12 月 1 日). “河外星云”（"Extragalactic nebulae"）. *The Astrophysical Journal*, 64: 321。  

Hubble, Edwin (1929 年 3 月 15 日). “河外星云的距离与径向速度关系”（"A Relation Between Distance and Radial Velocity Among Extra-Galactic Nebulae"）. *PNAS*, 15 (3): 168–173。  

Penzias, A. A.; Wilson, R. W. (1965 年 7 月 1 日). “在 4080 Mc/s 下测得的多余天线温度”（"A Measurement of Excess Antenna Temperature at 4080 Mc/s"）. *The Astrophysical Journal*, 142: 419–421。  

Boggess, N. W. 等 (1992 年 10 月 1 日). “COBE 任务：发射两年后的设计与性能”（"The COBE mission – Its design and performance two years after launch"）. *The Astrophysical Journal*, 397: 420–429。  

Parker, Barry R. (1993). 《大爆炸的辩护：突破与障碍》（*The Vindication of the Big Bang: Breakthroughs and Barriers*）. 纽约：Plenum Press。  

“计算机图形学成就奖”（"Computer Graphics Achievement Award"）. *ACM SIGGRAPH 2018 Awards*, 2018 年 8 月 12 日, 加拿大温哥华。  

Paraficz, D.; Hjorth, J.; Elíasdóttir, Á. (2009 年 5 月 1 日). “使用北欧光学望远镜监测的五个双类星体结果”（"Results of optical monitoring of 5 SDSS double QSOs"）. *Astronomy & Astrophysics*, 499 (2): 395–408。  

Guth, Alan. 在 Edge 基金会访谈中提出此观点（"Edge Foundation Interview"），存档于 2016 年 4 月 11 日。  

Sample, Ian (2014 年 6 月 4 日). “引力波探测结论被质疑后化为尘埃”（"Gravitational waves turn to dust after claims of flawed analysis"）. *The Guardian*。  

Cowen, Ron (2015 年 1 月 30 日). “引力波发现正式被否认”（"Gravitational waves discovery now officially dead"）. *Nature*.  

Overbye, Dennis (2014 年 12 月 1 日). “新图像细化了对早期宇宙的理解”（"New Images Refine View of Infant Universe"）. *The New York Times*.  

Charles Kahn (1994). 《阿那克西曼德与希腊宇宙学的起源》（*Anaximander and the Origins of Greek Cosmology*）. 哈克特出版社（Hackett Publishing）。  

Aristotle (1914). 《论天》（*De Mundo*）. 牛津大学出版社。  

Yavetz, Ido (1998 年 2 月). “论欧多克索斯的同心球体”（"On the Homocentric Spheres of Eudoxus"）. *Archive for History of Exact Sciences*, 52 (3): 222–225。  

Crowe, Michael (2001). 《从古代到哥白尼革命的世界理论》（*Theories of the World from Antiquity to the Copernican Revolution*）. Dover.  

Hirshfeld, Alan W. (2004). “阿里斯塔克斯的三角”（"The Triangles of Aristarchus"）. *The Mathematics Teacher*, 97 (4): 228–231。  

Eastwood, Bruce S. (2007). 《秩序诸天：加洛林文艺复兴时期的罗马天文学与宇宙学》（*Ordering the Heavens: Roman Astronomy and Cosmology in the Carolingian Renaissance*）. Brill, 第 238–239 页。  

Gilbert, William (1893). 《论磁体》（*De Magnete*），第六卷第三章。P. Fleury Mottelay 译本。纽约：Dover Publications。
\end{enumerate}


